% GENERATED, edit the French unit or the translation table
\mxunit{automatismes-ch09}{Cross-tabulating two qualitative variables}

\begin{fluencysheet}{Thu~11~Mar.~2027}{Cross-tabulating two qualitative variables}
\mxfq{Warm-up · Quick}{PRO}{Express $\dfrac{24}{60}$ as a percentage.}{$\pc{40}$.}{}
\mxfq{Core · Quick}{LEC}{In a class of $30$ pupils, $18$ are girls. What is the relative frequency of girls? of boys?}{$\dfrac{18}{30}=\pc{60}$; $\pc{40}$.}{}
\mxfq{Core + · Classify}{STA}{Classify each variable: eye colour; baccalauréat grade (AB, B, TB); height in cm.}{Nominal qualitative; ordinal qualitative (natural order); quantitative.}{}
\mxfq{Review · Problem}{MOD}{A rectangle of perimeter $20$ cm has length $x$. Express its area $A(x)$; calculate $A(3)$ and $A(5)$.}{Width $10-x$; $A(x)=x(10-x)$; $A(3)=21$, $A(5)=25$ cm$^2$.}{}
\mxfq{Thread · Quick}{SOL}{Calculate the exact volume of a cylinder of radius $2$ cm and height $5$ cm.}{$V=\pi\times4\times5=20\pi$ cm$^3$.}{}
\end{fluencysheet}

\begin{fluencysheet}{Mon~15~Mar.~2027}{Qualitative variables; two-way frequency table}
\mxfq{Warm-up · Quick}{CAL}{Complete: $12+\ldots=30$ and $22+\ldots=40$.}{$18$; $18$.}{}
\mxfq{Core · Table}{LEC}{Total $40$ pupils. Girls: $22$; canteen users (C): $25$; girls who use the canteen: $15$. Complete the two-way table.\par\smallskip\mxarraystretch{1.15}\begin{tabular}{|c|c|c|c|}\hline
 & C & Not C & Total \\ \hline
Girls & 15 & \ldots & 22 \\ \hline
Boys & \ldots & \ldots & \ldots \\ \hline
Total & 25 & \ldots & 40 \\ \hline
\end{tabular}\par }{Girls not using the canteen $7$; boys using the canteen $10$; boys not using the canteen $8$; total boys $18$; total not using the canteen $15$ (check: $7+8=15$).}{}
\mxfq{Core + · True/false}{LEC}{True or false? Justify. ``Satisfaction level (low, medium, high)'' is a nominal qualitative variable.}{False. There is a natural order: it is an ordinal qualitative variable.}{}
\mxfq{Review · Problem}{CMP}{Discount A: $-\pc{30}$ on $\num{80000}$ Ar. Discount B: $-\pc{20}$ then $-\pc{10}$. Which is the better deal?}{A: $\num{56000}$ Ar; B: $\num{80000}\times\num{0,8}\times\num{0,9}=\num{57600}$ Ar. A is the better deal.}{}
\mxfq{Thread · Method}{DEM}{We know that $AB^2+AC^2=BC^2$. What conclusion can be drawn? Which property is used?}{$ABC$ is right-angled at $A$: converse of Pythagoras' theorem.}{}
\end{fluencysheet}

\begin{fluencysheet}{Tue~16~Mar.~2027}{Marginal frequencies}
\mxfq{Warm-up · Quick}{PRO}{Express $\dfrac{25}{40}$ as a percentage.}{$\pc{62,5}$.}{}
\mxfq{Core · Method}{LEC}{Using the table from sheet 50: calculate the marginal frequency of girls and that of canteen users.}{Girls: $\dfrac{22}{40}=\pc{55}$; canteen: $\dfrac{25}{40}=\pc{62,5}$ (marginal counts over the total count).}{Marginal frequency = row (or column) total / overall total.}
\mxfq{Core + · Quick}{LEC}{The marginal frequency of girls is $\pc{55}$. What is that of boys?}{$\pc{45}$: the marginal frequencies of a variable add up to $\pc{100}$.}{}
\mxfq{Review · Quick}{SUB}{Calculate $(2a-b)^2$ for $a=3$ and $b=-2$.}{$(6+2)^2=64$.}{}
\mxfq{Thread · Quick}{REP}{$A\pt{-1;2}$ and $B\pt{3;2}$. What can be said about the line $(AB)$? Give $AB$.}{Horizontal (same ordinate), parallel to the horizontal axis; $AB=3-(-1)=4$.}{}
\end{fluencysheet}

\begin{fluencysheet}{Thu~18~Mar.~2027}{Conditional frequencies}
\mxfq{Warm-up · Quick}{PRO}{Give $\dfrac{15}{22}$ as a percentage rounded to the nearest tenth.}{$\pc{68,2}$.}{}
\mxfq{Core · Method}{LEC}{With the same table: among the girls, what is the frequency of canteen users? Among the boys?}{Girls: $\dfrac{15}{22}\approx\pc{68,2}$; boys: $\dfrac{10}{18}\approx\pc{55,6}$ (divide by the total of the subpopulation).}{}
\mxfq{Core + · Distinguish}{LEC}{Distinguish between ``the frequency of girls among canteen users'' and ``the frequency of canteen users among girls''.}{$\dfrac{15}{25}=\pc{60}$ and $\dfrac{15}{22}\approx\pc{68,2}$: the reference total is not the same.}{}
\mxfq{Review · Construction}{CON}{Draw $[AB]$ of length $6$ cm, then the circles with centres $A$ and $B$ and radius $4$ cm, which meet at $C$ and $D$. What type of quadrilateral is $ACBD$?}{A rhombus: $AC=CB=BD=DA=4$ cm (four equal sides).}{}
\mxfq{Thread · Quick}{DIV}{Is $91$ prime?}{No: $91=7\times13$.}{Pitfall: it is divisible by none of $2$, $3$, $5$.}
\end{fluencysheet}

\begin{fluencysheet}{Mon~22~Mar.~2027}{Completing a two-way table}
\mxfq{Warm-up · Quick}{PRO}{Calculate $\pc{35}$ of $80$.}{$28$.}{}
\mxfq{Core · Method}{LEC}{$200$ customers: $\pc{60}$ are women; $\pc{25}$ of the women and $\pc{50}$ of the men buy online. Complete the two-way table.}{Women $120$, of whom $30$ online; men $80$, of whom $40$ online; online in total: $70$.}{To go from conditional frequencies to counts, multiply by the size of the subpopulation.}
\mxfq{Core + · Check}{LEC}{A table states ``row total: $45+55=100$'' but ``column total: $60+30=90$''. What can be concluded?}{The table is inconsistent: the grand total must be the same by rows and by columns; there is a data-entry or calculation error.}{}
\mxfq{Review · True/false}{PRB}{True or false? Justify. After three consecutive ``heads'', the probability of getting ``tails'' on the next toss is greater than $\dfrac12$.}{False. The tosses are independent: the probability remains $\dfrac12$.}{}
\mxfq{Thread · Quick}{GEO}{$ABC$ is right-angled at $A$, $BC=10$ cm, $\widehat{ABC}=30\degre$. Calculate $AC$.}{$AC=BC\times\sin30\degre=10\times\num{0,5}=5$ cm.}{Use the sine: opposite side / hypotenuse.}
\end{fluencysheet}

\begin{fluencysheet}{Tue~23~Mar.~2027}{Real data: processing a file and filtering (AND, OR, NOT)}
\mxfq{Warm-up · Quick}{LEC}{In a file of $\num{1000}$ rows, $400$ satisfy $A$, $300$ satisfy $B$ and $100$ satisfy both $A$ and $B$. How many satisfy ``$A$ AND $B$''? ``$A$ OR $B$''?}{$100$; $400+300-100=600$ (the intersection must not be counted twice).}{}
\mxfq{Core · Method}{LEC}{With the same data: how many satisfy ``NOT $A$''? ``$A$ AND NOT $B$''?}{$\num{1000}-400=600$; $400-100=300$.}{}
\mxfq{Core + · Reading code}{ALG}{What is the final value of \code{compte}?\mxpseudo{compte = 0\\ for age in [15, 16, 17, 16, 18]:\\ \mxtab if age \textgreater{}= 16 and age \textless{} 18:\\ \mxtab\mxtab compte = compte + 1}}{$3$ (the ages $16$, $17$, $16$).}{}
\mxfq{Review · Quick}{NUM}{Give the decimal form of $\dfrac78$, then the fraction in lowest terms equal to $\num{0,12}$.}{$\num{0,875}$; $\dfrac{12}{100}=\dfrac{3}{25}$.}{}
\mxfq{Thread · Formula → table}{REG}{$f(x)=x^2-2x$. Complete a table of values for $x\in\ens{0;1;2;3}$.}{$0\ ;\ -1\ ;\ 0\ ;\ 3$.}{}
\end{fluencysheet}

\begin{fluencysheet}{Thu~25~Mar.~2027}{Critical reading of statements from the media}
\mxfq{Warm-up · Quick}{PRO}{A $\pc{10}$ increase on $40$: what is the final value?}{$40\times\num{1,10}=44$.}{}
\mxfq{Core · Critical thinking}{LEC}{An article claims: ``$\pc{60}$ of accidents involve sober drivers''. Can we conclude that driving sober is dangerous?}{No: the reference population is missing. If the vast majority of drivers are sober, it is normal for them to be the majority among those involved in accidents. Conditional frequencies would have to be compared.}{}
\mxfq{Core + · True/false}{LEC}{True or false? Justify. ``Two smokers out of three are men'' means that two men out of three smoke.}{False. This confuses $P_{\text{smoker}}(\text{man})$ and $P_{\text{man}}(\text{smoker})$: the reference populations differ.}{}
\mxfq{Review · True/false}{STA}{True or false? Justify. The median of a data set is always equal to its mean.}{False. Counter-example: $1,\,2,\,9$: mean $4$, median $2$.}{}
\mxfq{Thread · Reading code}{ALG}{What is the value of \code{s} at the end?\mxpseudo{s = 0\\ for i in range(1, 5):\\ \mxtab s = s + i}}{$s=1+2+3+4=10$.}{Pitfall: \code{range(1, 5)} stops at $4$; $5$ is excluded.}
\end{fluencysheet}
